\documentclass[10pt,a4paper]{amsart}
\usepackage[margin=19mm]{geometry}
\usepackage{amsmath,amssymb,mathtools}
\usepackage[scr=rsfs,cal=euler]{mathalfa}
\usepackage{xfrac}
\usepackage[unicode,hidelinks,bookmarks=false]{hyperref}
\newcommand{\ie}{i.e.\ }
\newcommand{\cf}{cf.\ }
\newcommand{\resp}{respectively\ }
\newcommand{\Subsection}{\subsection}
\providecommand{\nobreakdash}{\nobreak\hbox{-}}
\setlength{\emergencystretch}{2em}
\multlinegap0pt
\usepackage{amssymb}
\usepackage{enumerate}
\newtheorem{corollary}{Corollary}
\newtheorem{lemma}{Lemma}
\newtheorem{enumeratenumeric}{Enumeratenumeric}
\newcommand{\N}{\mathbb{N}}
\newcommand{\ad}{\operatorname{ad}}
\newcommand{\cD}{\mathcal{D}}
\newcommand{\oO}{\mathcal{O}}
\newcommand{\ord}{\mathrm{ord}}
\newcommand{\tmop}[1]{{\operatorname{#1}}}
\newcommand{\um}{{\underline m}}
\title{Bruno ideal and the variety of centers for singular germs of vector fields}
\author{Mar\'ia Mart\'in-Vega, Daniel Panazzolo}
\date{Source: arXiv:2602.13527v2 (CC BY 4.0)}
\begin{document}
\maketitle

\noindent\emph{Source and licence.} Bruno ideal and the variety of centers for singular germs of vector fields. Version arXiv:2602.13527v2 (CC BY 4.0), distributed by arXiv under the Creative Commons Attribution 4.0 International licence, http://creativecommons.org/licenses/by/4.0/, as recorded in the arXiv licence metadata for this exact version and shown on the abstract page https://arxiv.org/abs/2602.13527v2. Full licence notice: \texttt{licenses/arxiv-2602.13527-CC-BY-4.0.txt}.\par\smallskip

\noindent\textit{Editorial scope.} Counted unit: the article's lemma lemma-bound-expU (source0.tex lines 1331--1344). The article's complete original proof of it is delivered with it (source0.tex lines 1345--1397, one \texttt{proof} environment of 53 lines). No mathematics is written, repaired or re-derived: the statement and the proof are the article's own lines, in the article's own notation, and no retained line is substituted or re-flowed.  Retained uncounted as the source-local context the proof needs: lines 1254--1259; lines 1261--1263; lines 1295--1303.  Nothing was omitted from any retained span, so the package carries no omission record.  No retained line contains a citation, so the wrapper carries no bibliography and no bibliography entry.  No retained line contains a figure environment, an image-inclusion command or drawing code, so no image is delivered and the package declares no asset dependency.  Editorial formatting only, in addition: an AMS \texttt{amsart} preamble that restates the article's own declarations and macros used by the retained lines, namely the environments \texttt{corollary}, \texttt{lemma}, \texttt{enumeratenumeric} on separate counters, together with the standard packages those lines need.  The retained environments print in this document as Corollary 1, Lemma 1, Enumeratenumeric 1 in order of appearance, whereas the article prints the same items with the numbering of its own full document.  The complete native source of the article as distributed in the arXiv e-print for this exact version is bundled unchanged as \texttt{sources/194651/source0.tex}.
\begin{corollary}\label{corollary-inclusion-of-disks}
Let $\partial$ be as above and let $x \mapsto \varphi(x) = \Phi_\partial(1,x)$ denote the time-1 flow of $\partial$. Then
\[
\varphi(D_{r-c}) \subset D_r.
\]
\end{corollary}

\begin{equation}\label{estimate-f-flow}
\| \exp(\partial)(f) \|_{r-c} \le \| f \|_r, 
\end{equation}

\begin{corollary}\label{Corollary-estimate-deg-ord}
Let $\partial$ be a polynomial derivation. Then, for all $0<\rho<r$ and $n \in \N$, 
$$
\frac{1}{n!}\, \|  \ad_\partial^n (L (\mu))  \|_{\rho} \leq
\left[ \deg (\partial)  \left( \frac{\rho}{r}\right)^{\tmop{ord} (\partial)} \|\partial\|_r
\right]^n 
 \| \mu \|.
$$
\end{corollary}

\begin{lemma}\label{lemma-bound-expU}
  Under the above notation, there exists a $k_0 \in \mathbb{N}$ $($depending only on $C$ and the sequence $(\Omega_k)\, )$ such that,
  for all integers $k \geq k_0$, the following conditions hold:
  \begin{enumeratenumeric}
    \item $\varphi (D_{\rho_1}) \subset D_r$, where $\varphi$ is the time-1
    flow of $U$. 
    \item For an arbitrary function $f \in \oO (D_{\rho})$, \ \
    \[ \| \exp (U) f \|_{\rho_1} \leq \| f \|_{\rho}. \]
    \item If $R \in \cD (D_{\rho})$ is a derivation with $\ord (R) \geq
    1$ then
        \[ \| \exp (\tmop{ad}_U) R \|_{\rho_1} \leqslant \left( 1 - \frac{C k}{2^k}
   \right)  \| R \|_{\rho}. \]
 \end{enumeratenumeric}
\end{lemma}

\begin{proof}
  The item $(1)$ will follow directly from Corollary
  \ref{corollary-inclusion-of-disks} if we show that $\rho_1 \leq r - r
  \| U \|_r$. Under the above hypothesis on $\|U\|_r$, this is equivalent to prove the inequality
  \[ 1 - \left( \frac{1}{C 2^k} \right)^{C / 2^k} \geq \frac{1}{2^k}. \]
  But notice that the left-hand side of this inequality is equivalent to
  $\frac{k C\log 2}{2^k}$ as $k \rightarrow \infty$. Therefore, the
  inequality holds for all sufficiently large $k$.
  
  The item $(2)$ is an immediate consequence of item (1) and equation
  (\ref{estimate-f-flow}).
  
  Let us prove item (3). We initially consider the case of a monomial
  derivation $R = x^{\um} L (\mu)$ and observe that we can write
  \[ \exp (\ad_U) \left( x^{\um} L (\mu) \right) = \exp (U) x^{\um}
     \cdot \exp (\ad_U) L (\mu). \]
  Therefore, applying separately item (2) of the present Lemma to the monomial $x^{\um}$ and Corollary
  \ref{Corollary-estimate-deg-ord} to the derivation $L (\mu)$, we obtain
  \begin{eqnarray*}
    \left\| \exp (\ad_U) \left( x^{\um} L (\mu) \right)
    \right\|_{\rho_1} & \leq & \left\| \exp (U) x^{\um} \right\|_{\rho_1}
    \cdot \| \exp (\ad_U) L (\mu) \|_{\rho_1}\\
    & \leq & \left\| x^{\um} \right\|_r  \sum_{n \geq 0} \left[
    \deg (U)  \left( \frac{\rho_1}{r} \right)^{\tmop{ord} (U)}  \| U \|_r
    \right]^n \| \mu \|\\
    & \leq & \left\| x^{\um} \right\|_r \sum_{n \geq 0} \left[ 2
    \left( \frac{1}{C 2^k} \right)^C  \right]^n\, \| \mu \|.
  \end{eqnarray*}
  We can assume that the constant $A = 2 \left( \frac{1}{C 2^k} \right)^C$ is
  $< 1$ by choosing $k$ sufficiently large. Therefore, applying property (1)
  to $x^{\um}$ and summing the geometric series, we obtain the estimate
  \[ \left\| \exp (\ad_U) \left( x^{\um} L (\mu) \right)
     \right\|_{\rho_1} \leq \frac{\left( r/\rho \right)^{\left\|
     \um \right\|}}{1 - A}  \left\| x^{\um} \right\|_{\rho} \| \mu \| =
     \frac{\left( r/\rho \right)^{\left\| \um \right\|}}{1 - A} 
     \left\| x^{\um} L (\mu) \right\|_{\rho} . \]
  In the general case, under the hypothesis that $\ord(R) \geq 1$, we
  can write the monomial expansion of $R$ as
  \[ R = \sum_{\left\| \um \right\| \geq 1} x^{\um} L \left( \lambda_{\um}
     \right). \]
  Applying the above estimate to each monomial derivation $x^{\um} L \left(
  \lambda_{\um} \right)$, we conclude that
  \[ \| \exp (\ad_U) R \|_{\rho_1} \leq \frac{r / \rho}{1 - A}  \|
     R \|_{\rho}. \]
Using that $\frac{1}{1 - A} \leqslant 1 + 2 A$ for $A \leqslant 1 / 2$ we obtain the equivalence
\[ \frac{r / \rho}{1 - A} \leqslant \Omega_k \frac{1}{(C 2^k)^{C / 2^k}}
   \left( 1 + 4 \frac{1}{C^C 2^{k C}} \right) = 1 - \frac{C k \log 2}{2^k}(1+o(1)),
\]
where $o(1)$ indicates a term which goes to zero as $k \rightarrow \infty$. As a consequence, there exists a $k_0 \in \N$ such that
\[ \| \exp (\tmop{ad}_U) R \|_{\rho_1} \leqslant \left( 1 - \frac{C k}{2^k}
   \right)  \| R \|_{\rho}, \]
for all $k \ge k_0$.      
\end{proof}

\end{document}
